% Prozentrechnung · Prozentsatz · Prüfungs-Fokus MSA – bank/prozentrechnung/e2.jsonl (zusammenbau v0.6, Prüfungs-Fokus)
\einheitenkopf[e2]{Einheit 2 · Prozentsatz berechnen}
\zweigzeile{ab Kl. 7 · P10 ×5}
\setcounter{aufgabe}{0}

% Anlauf (ohne Prüfkennung)
\begin{aufgabe}{Prozentsatz – Anlauf}
\begin{teile}
\teil Streifen: $50$ Kinder, grau: $10$ Kinder – wie viel Prozent?

\streifenfeld{20}{$0$}{$50$ Kinder}
\teil $17$ von $100$ Schülern spielen ein Instrument – wie viel Prozent? \leerfeld[\%]
\teil $19$ von $50$ Kindern – wie viel Prozent? \leerfeld[\%]
\end{teile}
\end{aufgabe}

\begin{aufgabe}{Prozentsatz – Prüfungsaufgaben}
\begin{teile}
\teil Bei einer Tombola gibt es $21$ Nieten und $3$ Hauptgewinne. Wie viel Prozent der Lose sind Hauptgewinne? \hfill (P18) \\ \leerfeld[\%]
\teil Auf dem Blech liegen $11$ Muffins mit Blaubeeren und $4$ mit Kirschen. Wie viel Prozent sind Kirschmuffins? Runde auf eine Stelle. \hfill (P18) \\ \leerfeld[\%]
\teil Einwohner einer Stadt: gesamt $480\,000$; Kinder $72\,000$; Jugendliche $57\,000$; Erwachsene $264\,000$; Senioren $87\,000$. Wie viel Prozent sind Jugendliche? Runde auf eine Stelle. \hfill (P23)
\teil Wasser eines Haushalts pro Tag: gesamt $380$ l; Toilette $102$ l; Duschen $136$ l; Wäsche $46$ l; Rest $96$ l. Wie viel Prozent entfallen aufs Duschen? Runde auf eine Stelle. \hfill (P23)
\teil Von $64$ Bewerbungen wurden $47$ angenommen. Wie viel Prozent wurden abgelehnt? Runde auf eine Stelle. \hfill (P15)
\teil Ein Zug fuhr an $92$ Tagen, $79$-mal war er pünktlich. An wie viel Prozent der Tage war er verspätet? Runde auf eine Stelle. \hfill (P15)
\teil Eine Stadt hat $40\,000$ Einwohner, $1\,080$ sind in diesem Jahr zugezogen. Wie viel Prozent sind das? \hfill (P14) \\ \leerfeld[\%]
\teil Von $12\,500$ verkauften Konzertkarten wurden $175$ zurückgegeben. Wie viel Prozent sind das? \hfill (P14) \\ \leerfeld[\%]
\teil Körpergrößen der $13$ Spieler einer Mannschaft in cm: $172$, $185$, $169$, $190$, $178$, $181$, $166$, $188$, $175$, $183$, $170$, $192$, $177$. Wie viel Prozent sind größer als $180$ cm? Runde auf eine Stelle. \hfill (P25)
\teil Wartezeiten von $9$ Patienten in Minuten: $12$, $35$, $8$, $22$, $41$, $17$, $30$, $5$, $26$. Wie viel Prozent warteten länger als $20$ Minuten? Runde auf eine Stelle. \hfill (P25)
\end{teile}
\end{aufgabe}

\medskip
\uebersichtskasten{Prozentsatz: Teil geteilt durch Ganzes – das ist der Anteil, dann in Prozent schreiben. \\ 14 von 40: 14 : 40 = 0,35 = 35 \% \quad 3 von 8: 3 : 8 = 0,375 = 37,5 \%}
